\begin{tomtat}
	\begin{minipage}[t]{0.3\textwidth}
		\vspace{0pt}
		\begin{tomtat}
			\begin{center}
				\textbf{ĐỀ 07}
			\end{center}
			
		\end{tomtat}
	\end{minipage}
	\hfill
	\begin{minipage}[t]{0.5\textwidth}
		\vspace{0pt}
		\begin{tomtat}
			\begin{center}
				\includegraphics[scale=0.08]{images/7}
			\end{center}
		\end{tomtat}
	\end{minipage}
\end{tomtat}
\textbf{I. PHẦN TRẮC NGHIỆM}
\begin{ex}
	Phương trình $(x-2)(x+3)=0$ có nghiệm là:
	\choice{$x \in \{2;-3\}$}{$x \in \{-2;3\}$}{$x \in \{-2;-3\}$}{$x \in \{2;3\}$}
	\loigiai{
		Ta có phương trình tích $(x-2)(x+3)=0$.
		Một tích bằng 0 khi và chỉ khi ít nhất một trong các thừa số của nó bằng 0.
		Do đó, ta có hai trường hợp:
		\begin{itemize}
			\item $x-2=0 \Rightarrow x=2$
			\item $x+3=0 \Rightarrow x=-3$
		\end{itemize}
		Vậy, tập nghiệm của phương trình là $S = \{2; -3\}$.
		\\ \textbf{ĐÁP ÁN A}
	}
\end{ex}

\begin{ex}
	Nghiệm của bất phương trình $3x+12 \le 0$ là:
	\choice{$x \le 4$}{$x \ge 4$}{$x \le -4$}{$x \ge -4$}
	\loigiai{
		Ta giải bất phương trình:
		\begin{align*}
			3x+12 &\le 0 \\
			3x &\le -12 \\
			x &\le \frac{-12}{3} \\
			x &\le -4
		\end{align*}
		Vậy nghiệm của bất phương trình là $x \le -4$.
		\\ \textbf{ĐÁP ÁN C}
	}
\end{ex}

\begin{ex}
	Căn bậc ba của 27 là:
	\choice{3}{-3}{3 và -3}{9}
	\loigiai{
		Căn bậc ba của một số $a$ là một số $x$ sao cho $x^3 = a$.
		Ta cần tìm số $x$ sao cho $x^3 = 27$.
		Ta thấy $3^3 = 3 \times 3 \times 3 = 9 \times 3 = 27$.
		Vậy căn bậc ba của 27 là 3.
		\\ \textbf{ĐÁP ÁN A}
	}
\end{ex}

\begin{ex}
	Tính $\sqrt{\frac{4}{9}}$ ta được kết quả là:
	\choice{$\sqrt{\frac{2}{3}}$}{$-\sqrt{\frac{2}{3}}$}{$\frac{2}{3}$}{$\pm\frac{2}{3}$}
	\loigiai{
		Áp dụng quy tắc khai phương một thương, ta có:
		$$ \sqrt{\frac{4}{9}} = \frac{\sqrt{4}}{\sqrt{9}} $$
		Vì $\sqrt{4} = 2$ (do $2^2=4$ và $2>0$) và $\sqrt{9} = 3$ (do $3^2=9$ và $3>0$).
		Nên:
		$$ \sqrt{\frac{4}{9}} = \frac{2}{3} $$
		\\ \textbf{ĐÁP ÁN C}
	}
\end{ex}

\begin{ex}
	Trong các công thức sau. Công thức nào là công thức tính diện tích hình quạt tròn?
	\choice{$S=\pi R^2$}{$l=\frac{\pi Rn}{180}$}{$S=\frac{\pi R^2n}{360}$}{$S=2\pi(R^2-r^2)$}
	\loigiai{
		Công thức tính diện tích hình quạt tròn bán kính $R$, cung $n^\circ$ là:
		$S = \frac{\pi R^2 n}{360}$.
		Trong đó:
		\begin{itemize}
			\item A. $S=\pi R^2$ là công thức tính diện tích hình tròn.
			\item B. $l=\frac{\pi Rn}{180}$ là công thức tính độ dài cung tròn $n^\circ$.
			\item D. $S=2\pi(R^2-r^2)$ không phải là công thức chuẩn. Diện tích hình vành khuyên là $S=\pi(R^2-r^2)$.
		\end{itemize}
		Vậy công thức đúng là C.
		\\ \textbf{ĐÁP ÁN C}
	}
\end{ex}

\begin{ex}
	Gieo một con xúc sắc 50 lần và được kết quả như sau:
	\begin{center}
		\begin{tabular}{|l|c|c|c|c|c|c|}
			\hline
			Số chấm xuất hiện & 1 & 2 & 3 & 4 & 5 & 6 \\
			\hline
			Tần số & 8 & 7 & 9 & 8 & 6 & 12 \\
			\hline
		\end{tabular}
	\end{center}
	Tần số xuất hiện mặt năm chấm là:
	\choice{9}{98}{7}{6}
	\loigiai{
		Dựa vào bảng thống kê kết quả gieo xúc xắc:
		Khi "Số chấm xuất hiện" là 5, "Tần số" tương ứng là 6.
		Vậy tần số xuất hiện mặt năm chấm là 6.
		\\ \textbf{ĐÁP ÁN D}
	}
\end{ex}

\begin{ex}
	Lương của các công nhân một nhà máy được cho trong bảng sau:

		\begin{tabular}{|l|c|c|c|c|c|}
			\hline
			Lương & [5;7) & [7;9) & [9;11) & [11;13) \\
			\hline
			Tần số& 20 & 50 & 70 & 40 \\
			\hline
		\end{tabular}

	Để vẽ biểu đồ tần số tương đối ghép nhóm dạng đoạn thẳng, ta dùng giá trị nào đại diện cho nhóm số liệu [9;11)?
	\choice{11}{10,5}{10}{9}
	\loigiai{
		Để vẽ biểu đồ tần số tương đối ghép nhóm dạng đoạn thẳng, giá trị đại diện cho mỗi nhóm số liệu $[a;b)$ thường được chọn là trung điểm của khoảng đó, tức là $\frac{a+b}{2}$.
		Đối với nhóm số liệu $[9;11)$, giá trị đại diện là:
		$$ \frac{9+11}{2} = \frac{20}{2} = 10 $$
		\\ \textbf{ĐÁP ÁN C}
	}
\end{ex}

\begin{ex}
	Công thức tính diện tích xung quanh hình trụ là:
	\choice{$S = \frac{1}{2}Cl = \pi Rl$}{$V = \frac{1}{3}Sh = \frac{1}{3}\pi R^2 h$}{$S = 4\pi R^2$}{$S = Ch = 2\pi R h$}
	\loigiai{
		Công thức tính diện tích xung quanh của hình trụ có bán kính đáy $R$ và chiều cao $h$ là $S_{xq} = 2\pi Rh$.
		Trong đó $C = 2\pi R$ là chu vi đáy. Vậy $S_{xq} = C \cdot h$.
		Xét các lựa chọn:
		\begin{itemize}
			\item A. $S = \frac{1}{2}Cl = \pi Rl$ là công thức tính diện tích xung quanh hình nón (với $l$ là đường sinh).
			\item B. $V = \frac{1}{3}Sh = \frac{1}{3}\pi R^2 h$ là công thức tính thể tích hình nón.
			\item C. $S = 4\pi R^2$ là công thức tính diện tích mặt cầu.
			\item D. $S = Ch = 2\pi R h$ là công thức tính diện tích xung quanh hình trụ.
		\end{itemize}
		Vậy công thức đúng là D.
		\\ \textbf{ĐÁP ÁN D}
	}
\end{ex}
\textbf{II. PHẦN TỰ LUẬN}
\begin{ex}
	Giải hệ phương trình: $\begin{cases} 7x+4y=18 \\ 3x-4y=2 \end{cases}$
	\loigiai{
		Ta có hệ phương trình:
		$$ \begin{cases} 7x+4y=18 & (1) \\ 3x-4y=2 & (2) \end{cases} $$
		Cộng vế theo vế của phương trình (1) và phương trình (2), ta được:
		\begin{align*}
			(7x+4y) + (3x-4y) &= 18 + 2 \\
			10x &= 20 \\
			x &= \frac{20}{10} \\
			x &= 2
		\end{align*}
		Thay $x=2$ vào phương trình (2):
		\begin{align*}
			3(2) - 4y &= 2 \\
			6 - 4y &= 2 \\
			-4y &= 2 - 6 \\
			-4y &= -4 \\
			y &= 1
		\end{align*}
		Vậy nghiệm của hệ phương trình là $(x; y) = (2; 1)$.
	}
\end{ex}

\begin{ex}
	Cho $x>0$ và $x \ne 1, x \ne 9$. Rút gọn biểu thức $B = \left(\frac{x+25}{x-3\sqrt{x}} - \frac{10}{\sqrt{x}-3}\right) : \frac{\sqrt{x}-5}{\sqrt{x}}$
	\loigiai{
		Với $x>0$ và $x \ne 1, x \ne 9$, ta có:
		\begin{align*}
			B &= \left(\frac{x+25}{\sqrt{x}(\sqrt{x}-3)} - \frac{10\sqrt{x}}{\sqrt{x}(\sqrt{x}-3)}\right) : \frac{\sqrt{x}-5}{\sqrt{x}} \\
			&= \frac{x+25-10\sqrt{x}}{\sqrt{x}(\sqrt{x}-3)} \cdot \frac{\sqrt{x}}{\sqrt{x}-5} \\
			&= \frac{x-10\sqrt{x}+25}{\sqrt{x}(\sqrt{x}-3)} \cdot \frac{\sqrt{x}}{\sqrt{x}-5} \\
			&= \frac{(\sqrt{x}-5)^2}{\sqrt{x}(\sqrt{x}-3)} \cdot \frac{\sqrt{x}}{\sqrt{x}-5}
		\end{align*}
		Vì $x>0$, $\sqrt{x} \ne 0$. Vì $x \ne 25$ (điều kiện ẩn để $\sqrt{x}-5 \ne 0$ cho phép rút gọn), ta có thể rút gọn $\sqrt{x}$ và $\sqrt{x}-5$.
		\begin{align*}
			B &= \frac{\sqrt{x}-5}{\sqrt{x}-3}
		\end{align*}
		Vậy $B = \frac{\sqrt{x}-5}{\sqrt{x}-3}$.
	}
\end{ex}

\begin{ex}
	Biết rằng phương trình bậc hai $x^2-3x-5=0$ có hai nghiệm phân biệt $x_1, x_2$.
	Không giải phương trình, tính giá trị của biểu thức $C = \frac{x_1^2+3x_2}{x_1^2+x_2^2}$.
	\loigiai{
		Phương trình đã cho là $x^2-3x-5=0$.
		Đây là phương trình bậc hai có dạng $ax^2+bx+c=0$ với $a=1, b=-3, c=-5$.
		Ta tính biệt thức $\Delta = b^2-4ac = (-3)^2 - 4(1)(-5) = 9 + 20 = 29$.
		Vì $\Delta = 29 > 0$ nên phương trình có hai nghiệm phân biệt $x_1, x_2$.
		
		Theo hệ thức Vi-ét, ta có:
		$$ \begin{cases} 
			x_1+x_2 = \dfrac{-b}{a} = \dfrac{-(-3)}{1} = 3 \\ 
			x_1x_2 = \dfrac{c}{a} = \dfrac{-5}{1} = -5 
		\end{cases} $$
		Vì $x_1$ là nghiệm của phương trình $x^2-3x-5=0$, nên $x_1^2 - 3x_1 - 5 = 0 \Rightarrow x_1^2 = 3x_1 + 5$.
		Ta có $x_1^2+x_2^2 = (x_1+x_2)^2 - 2x_1x_2 = (3)^2 - 2(-5) = 9 + 10 = 19$.
		
		Xét tử số của $C$: $x_1^2+3x_2 = (3x_1+5) + 3x_2 = 3(x_1+x_2) + 5$.
		Thay $x_1+x_2 = 3$ vào tử số: $3(3) + 5 = 9 + 5 = 14$.
		
		Vậy, giá trị của biểu thức $C$ là:
		$$ C = \frac{14}{19} $$
	}
\end{ex}

\begin{ex}
	Có ba chiếc hộp. Hộp A chứa 2 tấm thẻ lần lượt ghi số 1 và 2. Hộp B chứa 3 tấm thẻ lần lượt ghi các số 1; 2 và 3. Hộp C chứa 4 quả cầu lần lượt ghi các số 1; 2; 3 và 4. Các tấm thẻ cùng loại, các quả cầu có cùng kích thước và khối lượng. Bạn An rút ngẫu nhiên đồng thời một tấm thẻ từ mỗi hộp A và B. Bạn Bình lấy ngẫu nhiên một quả cầu từ hộp C. Tính xác suất của biến cố D: "Tổng ba số ghi trên hai tấm thẻ và quả cầu là 6."
	\loigiai{
		Tập hợp tất cả các kết quả có thể xảy ra (không gian mẫu $\Omega$) là các bộ ba $(a,b,c)$.
		Các phần tử của không gian mẫu $\Omega$:\\
			(1,1,1), (1,1,2), (1,1,3), (1,1,4)
			(1,2,1), (1,2,2), (1,2,3), (1,2,4)
			(1,3,1), (1,3,2), (1,3,3), (1,3,4)
			(2,1,1), (2,1,2), (2,1,3), (2,1,4)
			(2,2,1), (2,2,2), (2,2,3), (2,2,4)
			(2,3,1), (2,3,2), (2,3,3), (2,3,4)
		
		Tập hợp các kết quả thuận lợi cho biến cố D: "Tổng ba số ghi trên hai tấm thẻ và quả cầu là 6" (tức là $a+b+c=6$). (1,1,4); (1,2,3);(1,3,2);(2,1,3);(2,2,2);(2,3,1)
		
		Xác suất của biến cố D là:
		$$ P(D) = \frac{n(D)}{n(\Omega)} = \frac{6}{24} = \frac{1}{4} $$
	}
\end{ex}

\begin{ex}
	Một đội xe dự định chở 75 tấn hàng để ủng hộ đồng bào miền trung, lúc sắp khởi hành nhận được ủng hộ thêm 5 tấn hàng và được bổ sung thêm 5 xe, do đó mỗi xe chở ít hơn dự định 1 tấn. Hỏi dự định có bao nhiêu xe?
	\loigiai{
		Gọi số xe dự định của đội là $x$ (xe).
		Điều kiện: $x \in \mathbb{N}^*$.
		
		Theo dự định, mỗi xe chở số tấn hàng là: $\frac{75}{x}$ (tấn/xe).
		
		Lúc sắp khởi hành:
		Tổng số hàng cần chở là $75 + 5 = 80$ tấn.
		Số xe thực tế là $x+5$ (xe).
		Thực tế, mỗi xe chở số tấn hàng là: $\frac{80}{x+5}$ (tấn/xe).
		
		Vì mỗi xe thực tế chở ít hơn dự định 1 tấn, nên ta có phương trình:
		$$ \frac{75}{x} - \frac{80}{x+5} = 1 $$
		Quy đồng mẫu số:
		\begin{align*}
			\frac{75(x+5) - 80x}{x(x+5)} &= \frac{x(x+5)}{x(x+5)} \\
			75x + 375 - 80x &= x^2 + 5x \\
			-5x + 375 &= x^2 + 5x \\
			x^2 + 5x + 5x - 375 &= 0 \\
			x^2 + 10x - 375 &= 0
		\end{align*}
		Tính biệt thức $\Delta = b^2 - 4ac = 10^2 - 4(1)(-375) = 100 + 1500 = 1600$.
		$\sqrt{\Delta} = \sqrt{1600} = 40$.
		Phương trình có hai nghiệm:
		$$ x_1 = \frac{-10 + 40}{2(1)} = \frac{30}{2} = 15 $$
		$$ x_2 = \frac{-10 - 40}{2(1)} = \frac{-50}{2} = -25 $$
		Vì $x \in \mathbb{N}^*$, nên ta nhận $x=15$.
		Kiểm tra điều kiện:
		Số tấn mỗi xe dự định chở: $\frac{75}{15} = 5$ tấn.
		Số tấn mỗi xe thực tế chở: $\frac{80}{15+5} = \frac{80}{20} = 4$ tấn.
		Chênh lệch: $5 - 4 = 1$ tấn (đúng).
		
		Vậy, dự định có 15 xe.
	}
\end{ex}


\begin{ex} % Câu 16
	Cho tam giác nhọn $ABC$ nội tiếp trong đường tròn tâm $O$ có $AB < AC$. Trên cung nhỏ $AC$ lấy điểm $M$ khác $A$ thỏa mãn $MA < MC$. Vẽ đường kính $MN$ của đường tròn $(O)$ và gọi $H, K$ lần lượt là hình chiếu vuông góc của $A$ trên $DB, DE$. Chứng minh rằng:
	\begin{enumerate}
		\item[a)] Bốn điểm $A, H, K, M$ cùng nằm trên một đường tròn.
		\item[b)] $AH \cdot AK = HB \cdot MK$. và Đường thẳng $HK$ cắt $AB$ tại trung điểm của $AB$. 
	\end{enumerate}
\end{ex}

\begin{ex} % Câu 17
	\begin{enumerate}
		\item[a)] Phần mái lá của một ngôi nhà có dạng hình nón (không có đáy) với đường kính đáy khoảng 12m và độ dài đường sinh khoảng 8,5m (Hình 26). Chi phí để làm phần mái đó là 250 000 đồng/m$^2$. Hỏi tổng chi phí để làm toàn bộ phần mái nhà đó là bao nhiêu đồng?
		\item[b)] Cho các số thực dương $x, y, z$ thỏa mãn điều kiện $xy + yz + zx = 3xyz$.
		Tìm giá trị nhỏ nhất của biểu thức $A = \frac{x^2}{z(z^2+x^2)} + \frac{y^2}{x(x^2+y^2)} + \frac{z^2}{y(y^2+z^2)}$.
	\end{enumerate}
	\loigiai{
		\begin{center}
		\begin{tikzpicture}[line join = round, line cap=round,>=stealth,font=\footnotesize,scale=1]
			% Def
			\def\r{3} % r=3cm
			
			% Coordinates
			\coordinate (O) at (0,0);
			\coordinate (A) at (130:\r); % A(110 độ, 2.1*r)
			\coordinate (B) at (-160:\r);
			\coordinate (C) at (-20:\r);
			\coordinate (M) at (70:\r); % Đổi M_point thành M
			
			\coordinate (N) at ($(O)!-1!(M)$); % O la trung diem cua M N
			\coordinate (K) at ($(M)!(A)!(N)$); % K la hinh chieu cua A len M N
			\coordinate (H) at ($(M)!(A)!(B)$); % H la hinh chieu cua A len M B
			
			\coordinate (E) at ($(A)!0.5!(B)$); % E la trung diem cua AB (NEW DEFINITION)
			
			% Draw
			\draw (O) circle (\r); % Ve duong tron (O,r)
			\draw (A) -- (B) -- (C) -- cycle; % Ve tam giac ABC
			\draw (M) -- (N); % Ve M N
			\draw (B) -- (M); % Ve B M
			\draw (E) -- (K); % Ve EK
			\draw (A) -- (H); % Ve AH
			\draw (A) -- (K); % Ve AK
			
			% Ky hieu goc vuong bang lenh \pic
			% Can \usetikzlibrary{angles, quotes} o preamble
			\pic[draw, angle radius=2.5mm]{right angle = A--H--M}; % Goc AHM vuong
			\pic[draw, angle radius=2.5mm]{right angle = A--K--N};     % Goc AKN vuong
			
			% Fill 
			\fill (O) circle (1.5pt);
			\fill (A) circle (1.5pt);
			\fill (B) circle (1.5pt);
			\fill (C) circle (1.5pt);
			\fill (M) circle (1.5pt);
			\fill (N) circle (1.5pt);
			\fill (K) circle (1.5pt);
			\fill (H) circle (1.5pt);
			\fill (E) circle (1.5pt);
			
			% Nodes	
			\node[below left] at (O) {$O$};
			\node[above] at (A) {$A$};
			\node[below left] at (B) {$B$};
			\node[below right] at (C) {$C$};
			\node[above right] at (M) {$M$}; 
			\node[below left] at (N) {$N$};    % Vi tri nhan co the can dieu chinh
			\node[above right] at (K) {$K$};   % Vi tri nhan co the can dieu chinh
			\node[below right] at (H) {$H$};   % Vi tri nhan co the can dieu chinh
			\node[above] at (E) {$E$};       % Vi tri nhan co the can dieu chinh
			% Ve va danh dau doan thẳng bằng nhau
		\end{tikzpicture}
		
		\end{center}
		}
\end{ex}

\setcounterref{ex}{0}